%============================================================
\chapter{Einleitung}
\label{chap:einleitung}
%============================================================

\section{Motivation und Projekthintergrund}
\label{sec:motivation}

Autonomes Fahren hat sich in den vergangenen Jahren von einem
akademischen Forschungsfeld zu einer treibenden Kraft der modernen
Automobilindustrie entwickelt.
Fahrzeuge müssen dabei in der Lage sein, ihre Umgebung wahrzunehmen,
relevante Hindernisse zu erkennen und in Echtzeit sichere Fahrtrajektorien
zu berechnen.

Das vorliegende Projekt -- \textbf{Mxck Follow the Gap (FTG)} --
entstand im Rahmen einer Lehrveranstaltung und nutzt die
\emph{F1TENTH}-Plattform \cite{okelly2020f1tenth}: ein
maßstabsverkleinertes (1:10) Hochgeschwindigkeits-Rennfahrzeug,
das als standardisierte Plattform für die Entwicklung und den
Vergleich autonomer Fahrstrategien eingesetzt wird.

\subsection*{Zielsetzung}

Das übergeordnete Ziel des Projekts war es, das Fahrzeug
vollständig autonom auf einem abgesteckten Kurs fahren zu lassen
und dabei alle auftretenden Hindernisse sicher zu umfahren.
Als Kerntechnologie wurde der sogenannte
\emph{Follow-the-Gap}-Algorithmus (FTG) \cite{sezer2012novel}
gewählt, da er mit verhältnismäßig geringem Berechnungsaufwand
auskommt und ausschließlich auf \emph{LiDAR}-Rohdaten basiert.

\subsection*{Projektausgang}

Trotz intensiver Entwicklungs- und Testphasen konnte das
angestrebte Ziel -- eine robuste und vollständige
Hindernissvermeidung unter realen Bedingungen -- innerhalb des
vorgegebenen Zeitrahmens nicht vollständig erreicht werden.
Die Dokumentation schildert dennoch den gesamten Arbeitsprozess,
identifiziert die aufgetretenen Schwierigkeiten und gibt Empfehlungen
für weiterführende Arbeiten.

%------------------------------------------------------------
\section{Verwandte Arbeiten}
\label{sec:related_work}
%------------------------------------------------------------

Der Follow-the-Gap-Algorithmus wurde ursprünglich von
\citeauthor{sezer2012novel}~\cite{sezer2012novel} für mobile Roboter
vorgeschlagen und seitdem in verschiedenen Varianten für
Fahrzeuganwendungen adaptiert.
Im F1TENTH-Kontext stellt der Algorithmus eine beliebte Basisstrategie
dar, die in zahlreichen universitären Wettbewerben eingesetzt und
weiterentwickelt wurde \cite{fulgentini2022comparative}.

Gegenüber rein reaktiven Verfahren wie dem \emph{Vector Field Histogram}
(VFH) \cite{ulrich1998vfh} zeichnet sich FTG durch seine direkte
Nutzung der Lücken (\enquote{gaps}) im LiDAR-Scan aus, anstatt
Potenzialfelder zu konstruieren.

%------------------------------------------------------------
\section{Aufbau der Dokumentation}
\label{sec:structure}
%------------------------------------------------------------

Die Dokumentation ist wie folgt gegliedert:

\begin{itemize}[leftmargin=*]
  \item \textbf{Kapitel~\ref{chap:systemaufbau}} beschreibt die
        Hardware- und Softwareplattform des Mxck-Fahrzeugs.
  \item \textbf{Kapitel~\ref{chap:algorithmus}} erläutert den
        Follow-the-Gap-Algorithmus und dessen theoretische Grundlagen.
  \item \textbf{Kapitel~\ref{chap:implementierung}} dokumentiert die
        konkrete ROS\,2-Implementierung, die Knotenarchitektur und
        die verwendeten Parameter.
  \item \textbf{Kapitel~\ref{chap:ergebnisse}} stellt die Testergebnisse
        und Beobachtungen aus den Versuchsläufen dar.
  \item \textbf{Kapitel~\ref{chap:fazit}} fasst das Projekt zusammen,
        benennt Ursachen für das Nichterreichen des Ziels und gibt
        Empfehlungen für künftige Arbeiten.
\end{itemize}
